\chapter{绪论}

\section{研究背景}

跨媒介职业转型已经成为东亚娱乐产业中的常见现象。传统杂志、电视剧、综艺节目与视频平台在曝光节奏、互动强度和危机扩散速度上存在显著差异，因此艺人的职业转型不再只是职业名称的变化，而是生产机制、平台规则与公众期待的重新匹配。

\section{案例选择}

佐佐木希的公开职业轨迹具有较强的演示价值。其职业阶段跨度长，能够覆盖传统媒体与新媒体两种生产逻辑；角色身份横跨模特、演员、综艺嘉宾与视频创作者，适合展示多次转型的结构化写法；相关事件均有公开报道，便于在不暴露保密材料的前提下组织完整示例。

\section{研究目的}

本示例的目标有两层。第一，基于公开案例概括跨媒介职业转型的阶段特征、关键触发因素与形象韧性机制。第二，用可公开分享的内容验证 `thesis-just-bachelor` 在封面、摘要、目录、正文、参考文献和致谢等环节上的排版稳定性。
